\maketitle

\begin{abstract}
Let $f_3(n)$ be the maximum number of pairs at distance one among $n$ points in
$\R^3$ whose mutual distances are at least one.  We prove, for every integer
$n\ge2$,
\[
 \boxed{f_3(n)<6n-2.0465\,n^{2/3}.}
\]
More precisely,
\[
 f_3(n)\le6n-c_Kn^{2/3},
 \qquad
 c_K=
 \left(\frac{18}{\pi^2}\right)^{1/3}
 \left(103-\frac{117\sqrt3}{2}\right)
 =2.046505556385112\ldots .
\]
The global finite-volume estimate is proved by an FCC torus-completion
argument.  A finite packing and its radius-$r$ outer-parallel union are embedded
without wrap-around in a sufficiently large flat torus whose period lattice is a
coarse sublattice of the FCC lattice.  A random translate of the FCC packing is
then retained only outside the outer-parallel union.  Translation averaging and
the Kepler theorem give, for every $r\ge2$,
\[
 n\operatorname{vol}(B^3)
 \le\frac{\pi}{\sqrt{18}}\operatorname{vol}(U_r).
\]
Locally, the union of covering caps on an enlarged sphere is a fixed geodesic
neighborhood of the pairwise interior-disjoint radius-$\pi/6$ contact caps.
Spherical-neighborhood isoperimetry and a convex-chord argument reduce the
optimal degreewise affine charge to degrees one and eleven.  An exact minimax
calculation selects
\[
 r_*=(40+22\sqrt3)/37.
\]
All new algebra and the clean decimal comparison are accompanied by exact
rational certificates, and the algebraic local-to-global spine is Lean-checked
from explicit geometric hypotheses.  No lattice or crystallization assumption is
imposed on the finite extremizer.
\end{abstract}

\section{Introduction}

Let
\[
 f_3(n)=\max\bigl\{\#\{\{x,y\}\subset X:|x-y|=1\}: |X|=n,
 |x-y|\ge1\text{ for }x\ne y\bigr\}.
\]
After multiplying all distances by two, $f_3(n)$ is the maximum contact number of
a packing of $n$ congruent unit balls.  The kissing number in dimension three
gives $f_3(n)\le6n$.  Erd\H{o}s asked for the correct surface-order subtraction
\cite{Erdos1975}.  Bezdek and Reid proved the explicit estimate
\begin{equation}\label{eq:BR}
 f_3(n)<6n-0.926n^{2/3}
\end{equation}
using enlarged unions, truncated Voronoi cells, Euclidean isoperimetry, and a
local spherical-cap density bound \cite{BezdekReid2013}.  See also
\cite{BezdekKhan2018}.

The proof below has the same global/local architecture but uses two different
ingredients.  The global input is a finite completion principle.  Given a finite
unit-ball packing and the union $U_r$ of its concentric radius-$r$ balls, with
$r\ge2$, place the cluster on a sufficiently large flat torus obtained from a
coarse sublattice of the FCC lattice.  Randomly translate the FCC center set and
retain precisely those filler centers lying outside $U_r$.  Their expected hard
ball volume is
\[
 \delta_3\bigl(\operatorname{vol}(T)-\operatorname{vol}(U_r)\bigr),
\]
where $\delta_3=\pi/\sqrt{18}$.  The original cluster together with a suitable
translate of the retained filler is a torus packing.  Lifting it periodically and
applying the Kepler theorem yields
\[
 n\operatorname{vol}(B^3)\le\delta_3\operatorname{vol}(U_r).
\]
This directly proves the finite outer-parallel estimate needed here.  It is
related to the truncated-density language of Bezdek and L\'angi
\cite{BezdekLangi2025}, but it does not require a product identity for separately
taken upper limits.

The local input is spherical-neighborhood isoperimetry.  Around a fixed ball,
contact directions are separated by at least $\pi/3$, so their radius-$\pi/6$
caps have disjoint interiors.  If the packing balls are enlarged to radius $r$,
a tangent neighbor covers a cap of angular radius $\alpha=\arccos(1/r)$.  The
covering-cap union is exactly the $(\alpha-\pi/6)$-neighborhood of the disjoint
contact-cap union.  A spherical cap minimizes neighborhood measure among sets of
given measure, and the resulting local exposed-area expression is convex in the
cosine of the equal-area cap radius.

\begin{theorem}[universal contact bound]\label{thm:main}
For every integer $n\ge2$,
\begin{equation}\label{eq:mainexact}
 f_3(n)\le6n-c_Kn^{2/3},
 \qquad
 c_K=
 \left(\frac{18}{\pi^2}\right)^{1/3}
 \left(103-\frac{117\sqrt3}{2}\right).
\end{equation}
Moreover,
\begin{equation}\label{eq:clean}
 c_K=2.046505556385112\ldots>\frac{4093}{2000},
\end{equation}
and therefore
\[
 f_3(n)<6n-2.0465\,n^{2/3}.
\]
\end{theorem}

The radius is not chosen numerically.

\begin{theorem}[optimality within the one-radius method]\label{thm:optimality}
Within the one-radius, equal-area-cap, degreewise-affine method defined in
\cref{sec:method}, over every $r\ge2$, the coefficient is uniquely maximized at
\begin{equation}\label{eq:rstar}
 r_*=
 \frac{40+22\sqrt3}{37}=2.110949128824737\ldots .
\end{equation}
\end{theorem}

The optimality assertion is intentionally restricted to this method class.  No
claim is made that $c_K$ is the optimal universal contact-number coefficient, that
the normalized deficit has a limit, or that any close-packed lattice is globally
extremal for the finite contact problem.

\section{Normalization and connected maximizers}

Scale an admissible minimum-distance-one configuration by two and center a closed
unit ball at each scaled point.  Write the centers as $c_1,\ldots,c_n$.  Thus
$|c_i-c_j|\ge2$, with contact exactly at equality.  Let $E$ be the set of
undirected contacts, let $d_i$ be the contact degree at $c_i$, and put
\begin{equation}\label{eq:deficit}
 D=6n-|E|.
\end{equation}
The handshaking identity gives
\begin{equation}\label{eq:handshake}
 \sum_{i=1}^n(12-d_i)=12n-2|E|=2D.
\end{equation}

For fixed $n$, the attainable contact counts form a nonempty subset of the finite
set $\{0,\ldots,\binom n2\}$, so the maximum is attained.

\begin{lemma}[connected maximizer]\label{lem:connected}
Every contact-maximizing packing with $n\ge2$ has connected contact graph.
\end{lemma}

\begin{proof}
Suppose the graph is disconnected and partition its centers into two nonempty
unions of components, $A$ and $C$.  Every cross-distance is then strictly larger
than two.  Keep $A$ fixed and translate $C$ continuously in a direction that
brings a chosen cross-pair together.  At the first time at which the minimum
cross-distance reaches two, all cross-distances are still at least two, all old
contacts are preserved, and a new contact has appeared.  This contradicts
maximality.
\end{proof}

Hence, for a maximizing packing,
\begin{equation}\label{eq:degreerange}
 1\le d_i\le12,
\end{equation}
where the upper bound is the three-dimensional kissing-number theorem
\cite{SchuetteWaerden1953}.  This also covers $n=2$.
